\chapter{Cross-cutting Themes}
\label{ch:themes}

Six ideas recurred across the week, and in three cases two speakers reached the same territory from
opposite directions or disagreed outright. A disagreement between people who both know the field
marks an open question rather than a gap in the notes, so those three are the ones to dwell on.

\section*{1. The i.i.d.\ assumption, and a real disagreement about how bad it is}

Chandar closes the opening lecture (Chapter~\ref{ch:01-intro-deep-learning}) by naming the frontier
without explaining it: current practice optimises the i.i.d.\ case, while true learning is
sequential, online and continual. Three later lectures pick that up, and they do not agree.

Lyle (Chapter~\ref{ch:06-continual-learning}) treats it as an engineering fault with an identifiable
mechanism. Plasticity is lost because Adam's second-moment estimate is stale after a task change, so
the effective step size explodes, and because a large step pushes preactivations outside the
receptive field of the nonlinearity. Both are fixable by normalising, controlling parameter
norm, and avoiding regression on ill-conditioned targets, with resetting and distilling as a
fallback. Her verdict on forgetting is even more deflationary: if the data can be kept, the problem
is easy.

Pascanu (Chapter~\ref{ch:07-to-be-figured-out}) makes a structural claim instead. Gradient-based
credit assignment works by tug-of-war dynamics between competing examples, and that mechanism
\emph{requires} shuffled data. On this reading catastrophic forgetting is not a side effect but is
what failed credit assignment looks like, and Lyle's fixes are palliative because they stabilise
training without giving the learner explicit knowledge composition. He offers as hypotheses that
there is no compositional knowledge in these models and that scale is a crutch making learning
possible at all.

Deac's tutorial material (Chapter~\ref{ch:04-intro-reinforcement-learning}) is the third data point
and it supports Pascanu quietly: experience replay exists precisely because consecutive transitions
are correlated and SGD assumes independence. Replay is continual learning solved by refusing to be
in the continual setting, which is exactly Lyle's advice, and it is also an admission that nobody
knows how to learn from a correlated stream directly.

The disagreement is testable, which is what makes it interesting;
Chapter~\ref{ch:research-directions} sketches an experiment that would separate the two positions.

\section*{2. Symmetry imposed versus symmetry measured}

Veli\v{c}kovi\'{c} (Chapter~\ref{ch:11-geometric-deep-learning}) and Isola
(Chapter~\ref{ch:02-representational-geometry}) are doing the same thing from opposite ends.
Veli\v{c}kovi\'{c} specifies the symmetry group first and derives which layers are admissible,
strongly enough that CNNs are shown to be the only linear translation-equivariant maps on a grid.
Isola takes trained networks and measures, after the fact, which transformations their
representations are invariant to, then argues the right equivalence is isometry because distance is
what downstream tasks consume. One prescribes the invariance, the other discovers it.

The most striking evidence that this is one idea rather than two is that Cho
(Chapter~\ref{ch:05-statistics-estimation}) derives a permutation-equivariant architecture, a
transformer with positional embeddings removed, from exactly Veli\v{c}kovi\'{c}'s reasoning, in a
lecture about statistics, without either speaker referencing the other. The same argument
arriving independently in three lectures is some evidence it is not incidental.

\section*{3. Where knowledge gets injected}

Three lectures address the same question: if a model cannot learn everything from data, where does
the rest come from?

Cho's answer is to stop modelling the generative process. Inference need not invert generation, so an
estimator can be trained directly on problems with known answers, at the price of making
identifiability relative to the prior over problems, the hypothesis space and the optimiser. Pascanu's
answer is the mirror image: there is no prior-free machine learning system, every construction step
is a bias, so the question is never whether to encode knowledge but which encoding pays. And
Tapu\v{s}kovi\'{c} (Chapter~\ref{ch:08-ml-for-mathematics}) locates it in the evaluator, arguing that
in the 67-problem study search was never the bottleneck and that building a cheap, exact,
ungameable score \emph{is} the research.

Isola sits awkwardly across this. His lecture uses fixed metrics such as CKA to establish that
models converge, but his stated opinion is that there is no future for distance functions that are
not learned. Taken together with Cho, that points somewhere uncomfortable: if both the model and
the yardstick are fitted, convergence results become hard to state without circularity.

\section*{4. The spectrum as the design surface}

Three unrelated lectures turn out to be doing spectral engineering. Dieleman
(Chapter~\ref{ch:03-diffusion-models}) explains why diffusion works on images by observing that
natural images have a decaying power spectrum while Gaussian noise is flat, so the noise schedule is
effectively a frequency schedule and sampling recovers bands coarsest-first. Stankovi\'{c}
(Chapter~\ref{ch:09-spectral-analysis-dags}) cannot filter on a DAG at all until a spectrum is
restored, because nilpotency collapses every eigenvalue to zero. And Pascanu's linear recurrent unit
is designed entirely in the spectral domain, with the eigenvalue modulus setting the memory horizon
and the initialisation annulus chosen so that no unit begins with a memory too short to matter.

In all three the eigenvalues are not a diagnostic applied afterwards; they are the thing being
chosen.

\section*{5. Arithmetic is cheap, moving data is not}

Alistarh (Chapter~\ref{ch:10-efficient-learning}) and Pascanu converge on one conclusion from
opposite directions. Alistarh's framing is a 2500-fold gap between compute demand and hardware
capability over five years, and his remedy is to reduce bit-width so that weights occupy less of the
scarce resource. Pascanu profiles a TPU, finds roughly an eightfold gap between HBM and VMEM
transfer rates, and concludes the recurrent block should be cheap because the memory traffic
dominates. Neither is optimising FLOPs. Both treat arithmetic as abundant and data movement as the
constraint. The first question in any efficiency work is what has to move, not how many operations
there are.

\section*{6. Transformers are not what we call them}

Three lectures reframe the same architecture, and the reframings are compatible.
Veli\v{c}kovi\'{c} argues transformers are graph neural networks and not sequence models, with tokens
as nodes and the causal mask as a neighbourhood, which makes overmixing and oversquashing inherited
graph pathologies rather than novel LLM problems. Stankovi\'{c} notes in passing that a causal
attention mask is a strictly upper triangular matrix, hence a DAG, hence spectrally degenerate.
Pascanu argues transformers are a factorisation into simple components that must be scaled, and that
attention might not be all we need, while also reporting honestly that the hybrid alternative
retains 25\% global softmax attention because pure recurrence loses on language.

Read together they say the architecture is better understood as a message-passing operator on a
particular graph than as a sequence model, and that its known limitations follow from that graph.

\begin{aside}
One theme is conspicuous by its absence. Nobody argued for scale as a sufficient answer. Alistarh
opens with the bitter lesson and immediately derives a mandate for efficiency from it; Pascanu says
scale might work but is not the only option; Tapu\v{s}kovi\'{c} reports that the tool beats the
literature only where the literature is thin. That is a summer school's selection effect as much as a fact
about the field, but it is consistent across eleven speakers.
\end{aside}
